\section{Original coefficient compatibility of affine pullbacks into the function field}
\label{sec:scheme-function-field-pullback-coefficients}

\begin{theorem}[Literal native scalar compatibility under the original scheme morphism]
\label{thm:scheme-function-field-pullback-coefficients}
\lean{EnvelopeProperty.SchemeFunctionFieldPullbackCoefficients.sectionMap, EnvelopeProperty.SchemeFunctionFieldPullbackCoefficients.genericMap, EnvelopeProperty.SchemeFunctionFieldPullbackCoefficients.genericMap_apply}
\leanok
\uses{def:sc0-scheme-sections-function-field-map,thm:affine-function-field-generation}
\source{bj22-trivval:sec:tc}
For the actual morphism $f:X\to Y$ over the original field,
with original coefficient map $q:Y\to\operatorname{Spec}k$,
the unchanged pullback on sections over every original open
$V\subset Y$ is a $k$-algebra homomorphism for the original
coefficient structures induced by $q$ and $q\circ f$.
If $X$ is integral and $f^{-1}V$ is nonempty, composing this
same pullback with the actual generic germ gives a
$k$-algebra homomorphism into the same native $K(X)$.
Its value on every section is literally that original
pullback followed by that original generic germ.
\end{theorem}

\begin{proof}
The coefficient map on $\Gamma(Y,V)$ is the original map
$k\simeq\Gamma(\operatorname{Spec}k,\mathcal O)\to\Gamma(Y,V)$.
Compose it with the actual pullback map to
$\Gamma(X,f^{-1}V)$. The native composition identity for
pullback on open sections identifies this composite with
the coefficient map for the literal composite scheme
morphism $q\circ f$. Thus the unchanged ring homomorphism
commutes with the original $k$-actions and is a
$k$-algebra homomorphism. The actual generic germ is already
a $k$-algebra homomorphism for this same composite structure
map. Composing gives the asserted native field map, whose
pointwise formula is its definition. No independent
coefficient or function-field comparison is supplied.
\end{proof}
